\documentclass[10pt,a4paper]{amsart}
\usepackage[margin=19mm]{geometry}
\usepackage{amsmath,amssymb,mathtools}
\usepackage[scr=rsfs,cal=euler]{mathalfa}
\usepackage{xfrac}
\usepackage[unicode,hidelinks,bookmarks=false]{hyperref}
\newcommand{\ie}{i.e.\ }
\newcommand{\cf}{cf.\ }
\newcommand{\resp}{respectively\ }
\newcommand{\Subsection}{\subsection}
\providecommand{\nobreakdash}{\nobreak\hbox{-}}
\setlength{\emergencystretch}{2em}
\multlinegap0pt
\usepackage{amssymb}
\usepackage{mathtools}
\usepackage{tcolorbox}
\newtheorem{lemma}{Lemma}
\title{Powerfully embedded subgroups of extensions of powerful pro-$p$ groups and its applications to automorphisms of finite groups}
\author{Sathasivam Kalithasan, Tony N. Mavely, Viji Z. Thomas}
\date{Source: arXiv:2503.15240v3 (CC BY 4.0)}
\begin{document}
\maketitle

\noindent\emph{Source and licence.} Powerfully embedded subgroups of extensions of powerful pro-$p$ groups and its applications to automorphisms of finite groups. Version arXiv:2503.15240v3 (CC BY 4.0), distributed by arXiv under the Creative Commons Attribution 4.0 International licence, http://creativecommons.org/licenses/by/4.0/, as recorded in the arXiv licence metadata for this exact version and shown on the abstract page https://arxiv.org/abs/2503.15240v3. Full licence notice: \texttt{licenses/arxiv-2503.15240-CC-BY-4.0.txt}.\par\smallskip

\noindent\textit{Editorial scope.} Counted unit: the article's lemma Aut\_Phi(G) (source0.tex lines 1351--1353). The article's complete original proof of it is delivered with it (source0.tex lines 1355--1381, one \texttt{proof} environment of 27 lines). No mathematics is written, repaired or re-derived: the statement and the proof are the article's own lines, in the article's own notation, and no retained line is substituted or re-flowed.  No further source-local context is needed: every cross-reference of every retained line resolves inside the two delivered spans.  Nothing was omitted from any retained span, so the package carries no omission record.  No retained line contains a citation, so the wrapper carries no bibliography and no bibliography entry.  No retained line contains a figure environment, an image-inclusion command or drawing code, so no image is delivered and the package declares no asset dependency.  Editorial formatting only, in addition: an AMS \texttt{amsart} preamble that restates the article's own declarations and macros used by the retained lines, namely the environments \texttt{lemma} on separate counters, together with the standard packages those lines need.  The retained environments print in this document as Lemma 1 in order of appearance, whereas the article prints the same items with the numbering of its own full document.  The complete native source of the article as distributed in the arXiv e-print for this exact version is bundled unchanged as \texttt{sources/175048/source0.tex}.
\begin{lemma} \label{InnG powerfully embedded in Aut_Phi(G)}
Let $G$ be a finite $p$-group. If $G$ is powerful, then $\operatorname{Inn}(G)$ is powerfully embedded in $\operatorname{Aut}_\Phi(G)$. 
\end{lemma}

\begin{proof}
Let $g \in G$, and let $f_g$ be the inner automorphism defined by $g$. Note that $[f_g,\, \alpha]= f_{g\alpha(g^{-1})}$, for all $f_g \in \operatorname{Inn}(G)$ and $\alpha \in \operatorname{Aut}_\Phi(G)$.
%\[
%[f_g,\, \alpha] 
%= f_g \alpha f_g^{-1}\alpha^{-1} 
%= f_{g\alpha(g^{-1})}.
%\]
%Indeed, for any $x \in G$,
%\begin{align*}
%f_g \alpha f_g^{-1} \alpha^{-1}(x)
%&= f_g\bigl(\alpha(g^{-1})\, x\, \alpha(g)\bigr) \\
%&= g\alpha(g^{-1})\, x\, \alpha(g) g^{-1} \\
%&= g\alpha(g)^{-1}\, x\, \alpha(g) g^{-1} \\
%&= g\alpha(g^{-1}) \cdot x \cdot \bigl(g\alpha(g^{-1})\bigr)^{-1} \\
%&= f_{g\alpha(g^{-1})}(x).
%\end{align*}
Hence ${[f_g, \alpha] \in \operatorname{Inn}_{\Phi(G)}(G)}$. 
Since $\phi(G) = G^p$, it follows that
\[
[\operatorname{Inn} G,\, \operatorname{Aut}_\Phi(G)] \leqslant \operatorname{Inn}_{G^p}(G).
\]
We claim that $\operatorname{Inn}_{G^p}(G) 
\subseteq \operatorname{Inn}(G)^p$. Towards that end, let $h \in G^p$. Since $G$ is powerful, we have that $h=x^p$ for some $x\in G$. Thus ${f_h(g)=f_{x^p}(g)=(f_{x})^p(g)}$. Therefore
\[
[\operatorname{Inn} G,\, \operatorname{Aut}_\phi(G)] \leqslant \operatorname{Inn}(G)^p.
\]
\end{proof} 

\end{document}
